
This paper reports on a validation attempt for the hypothesis that incremental SMT verification would reduce verification time over batch re-verification for LLM code repair in typed Python with Pydantic. The validation pipeline failed before generating a single line of code due to API authentication errors, revealing that infrastructure robustness is a prerequisite to hypothesis testing.

The hypothesis---that incremental SMT verification (re-verifying only modified functions and dependencies) reduces verification time compared to batch re-verification---remains theoretically plausible but empirically unvalidated. All experimental validation was blocked at the foundational hypothesis (h-e1: static analyzers can extract SMT constraints from LLM-generated typed code) due to 100\% API authentication failure rate (48/48 generation attempts returned HTTP 401 ''API key is invalid'').

Despite generation failure, the dataset extension pipeline succeeded: HumanEval was mechanically extended with 100 Pydantic-annotated prompts. This artifact addresses a gap in formal verification research---prior benchmarks lack type annotations required for SMT constraint extraction. The constraint extraction pipeline (DatasetExtender $\rightarrow$ LLMGenerator $\rightarrow$ PyreExtractor $\rightarrow$ Z3Validator $\rightarrow$ MetricsEngine) was validated on error files, confirming architectural soundness on negative cases.

The contributions are:

\begin{enumerate}
\item \textbf{HumanEval + Pydantic Type Extensions}: 100 mechanically-generated type-annotated prompts suitable for SMT constraint extraction research, addressing the gap in typed benchmarks for formal verification of LLM code.
\end{enumerate}

\begin{enumerate}
\item \textbf{Constraint Extraction Pipeline Architecture}: A modular pipeline ready for retry, with components independently validated on error cases.
\end{enumerate}

The primary methodological contribution is the demonstration that infrastructure reliability is a prerequisite to hypothesis testing. Experimental pipelines must address failure modes explicitly through pre-flight API key validation, fast-fail on auth errors, and checkpoint/resume for long-running experiments.

\textbf{Honest Assessment of Claims}:
\begin{itemize}
\item Cannot claim ''incremental SMT achieves speedup'' (untested)
\item Cannot claim ''speedup scales with codebase size'' (untested)
\item Cannot claim ''conservative dependency analysis maintains soundness'' (untested)
\item Cannot claim ''type annotations enable constraint extraction from LLM code'' (untested)
\end{itemize}

All quantitative predictions await validation pending infrastructure fixes and completion of the h-e1 $\rightarrow$ h-m1 $\rightarrow$ h-m2 $\rightarrow$ h-m3 hypothesis chain.

\textbf{Limitations}: Single language scope (only Python + Pydantic + Pyre tested). Untested assumptions (A1: LLM code has extractable constraints; A2: repair locality transfers from human to LLM code). No baseline comparison (speedup figure is speculative, not measured).

\textbf{Future Work}: Retry h-e1 with valid Anthropic API key. If extraction\_rate $\geq 90$\%, proceed to mechanism testing. If extraction\_rate <90\%, pivot to neural constraint extraction. Extend approach to Rust + Prusti for language generalization. Validate repair locality assumption on LLM code. Standardize typed benchmarks for formal verification research across multiple languages.

The HumanEval + Pydantic dataset extension is the artifact from this validation attempt---a resource for future constraint extraction research, and a reminder that experimental methodology extends beyond hypothesis design to encompass infrastructure reliability.
